\section{Experimental Results}
Our experimental cohort contains 51 Transformer screening runs (17 configurations $\times$ 3 languages), 18 classical/dummy screening runs (6 $\times$ 3), 15 Transformer final runs (5 $\times$ 3), and 15 classical final runs (5 $\times$ 3). All use base seed 42. Tables are generated from stored metrics; fold summaries and per-label counts were checked arithmetically. Aggregate F1 uses four decimal places, and reported percentage-point differences use those displayed values. Per-label tables use three decimals. Section~\ref{sec:validity} gives the common evaluation qualifications.

\subsection{Traditional ML versus the Transformer reference: RQ1}
Table~\ref{tab:mlcv} and Figure~\ref{fig:mlvs} compare the shared Stage 1 folds. The strongest traditional method by mean micro-F1 is Random Forest for Python (0.5882), logistic regression for Java (0.5972), and SGD for Pharo (0.6486). E1 reaches 0.6169, 0.5968, and 0.7044, respectively. Relative to each language's strongest traditional model, E1 differs by approximately $+2.87$, $-0.04$, and $+5.58$ percentage points. These are descriptive differences between fold means, not significant effects.
\begin{table}[htbp]
\centering\small
\caption{Stage 1 micro-F1: mean and sample standard deviation across five folds. E1 uses fold validation for checkpoint selection.}\label{tab:mlcv}
\begin{tabular}{lrrr}\toprule
Method & Python & Java & Pharo\\\midrule
Logistic regression & 0.5833 $\pm$ 0.0126 & 0.5972 $\pm$ 0.0566 & 0.6438 $\pm$ 0.0223\\
Classifier chain & 0.5623 $\pm$ 0.0193 & 0.5789 $\pm$ 0.0675 & 0.6209 $\pm$ 0.0247\\
Linear SVM & 0.5730 $\pm$ 0.0095 & 0.5811 $\pm$ 0.0577 & 0.6372 $\pm$ 0.0210\\
SGD & 0.5811 $\pm$ 0.0093 & 0.5932 $\pm$ 0.0556 & 0.6486 $\pm$ 0.0188\\
Random forest & 0.5882 $\pm$ 0.0126 & 0.5726 $\pm$ 0.0405 & 0.6458 $\pm$ 0.0226\\
Majority dummy & 0.0000 $\pm$ 0.0000 & 0.0000 $\pm$ 0.0000 & 0.0000 $\pm$ 0.0000\\
RoBERTa E1 & 0.6169 $\pm$ 0.0177 & 0.5968 $\pm$ 0.0688 & 0.7044 $\pm$ 0.0374\\
\bottomrule\end{tabular}
\end{table}

\reportfigure{ml_vs_e1}{Five traditional methods versus RoBERTa E1 on Stage 1 micro-F1. Error bars are sample SD across five folds, not confidence intervals. LR denotes logistic regression.}{fig:mlvs}{1}

The Java logistic baseline also has macro-F1 0.6390 compared with 0.6273 for E1. The difference suggests that the more expensive reference encoder does not dominate lexical models under every metric. Pharo shows a clearer aggregate separation: E1 macro-F1 is 0.6053 versus 0.5604 for SGD, though this still does not prove better performance on every label. Full multi-metric screening tables are included in Appendix~\ref{app:metrics}.

\textbf{Answer to RQ1.} Traditional TF-IDF models are competitive with E1 on Java and provide meaningful baselines elsewhere. E1 has higher observed screening micro-F1 than all five traditional methods on Python and Pharo.

\subsection{Transformer configuration comparison: RQ2}
Table~\ref{tab:dlcv} covers every E1--E17 screening configuration. E15 (CodeBERT R-Drop) leads Python at 0.6479 and Java at 0.6225. E3 (RoBERTa focal) leads Pharo at 0.7152. Their improvements over E1 are about 3.10, 2.57, and 1.08 percentage points. There is no universal language-independent winner.
\begin{table}[htbp]
\centering\small
\caption{All 51 verified Transformer screening runs. E4 uses the same fixed threshold as E1 in Stage 1.}\label{tab:dlcv}
\begin{tabular}{lrrr}\toprule
Configuration & Python & Java & Pharo\\\midrule
E1 & 0.6169 $\pm$ 0.0177 & 0.5968 $\pm$ 0.0688 & 0.7044 $\pm$ 0.0374\\
E2 & 0.6161 $\pm$ 0.0199 & 0.5981 $\pm$ 0.0593 & 0.6912 $\pm$ 0.0445\\
E3 & 0.6219 $\pm$ 0.0256 & 0.5954 $\pm$ 0.0590 & 0.7152 $\pm$ 0.0233\\
E4 & 0.6169 $\pm$ 0.0177 & 0.5968 $\pm$ 0.0688 & 0.7044 $\pm$ 0.0374\\
E5 & 0.6080 $\pm$ 0.0213 & 0.6055 $\pm$ 0.0488 & 0.6992 $\pm$ 0.0318\\
E6 & 0.6158 $\pm$ 0.0167 & 0.6115 $\pm$ 0.0688 & 0.7036 $\pm$ 0.0278\\
E7 & 0.6081 $\pm$ 0.0102 & 0.5918 $\pm$ 0.0564 & 0.6736 $\pm$ 0.0208\\
E8 & 0.6271 $\pm$ 0.0285 & 0.6065 $\pm$ 0.0521 & 0.7104 $\pm$ 0.0307\\
E9 & 0.6211 $\pm$ 0.0268 & 0.6125 $\pm$ 0.0555 & 0.7124 $\pm$ 0.0304\\
E10 & 0.6260 $\pm$ 0.0196 & 0.5957 $\pm$ 0.0462 & 0.6950 $\pm$ 0.0253\\
E11 & 0.6397 $\pm$ 0.0181 & 0.5980 $\pm$ 0.0300 & 0.6805 $\pm$ 0.0358\\
E12 & 0.6132 $\pm$ 0.0181 & 0.6117 $\pm$ 0.0519 & 0.7068 $\pm$ 0.0367\\
E13 & 0.6455 $\pm$ 0.0163 & 0.6001 $\pm$ 0.0284 & 0.6893 $\pm$ 0.0347\\
E14 & 0.6307 $\pm$ 0.0196 & 0.6032 $\pm$ 0.0205 & 0.6767 $\pm$ 0.0282\\
E15 & 0.6479 $\pm$ 0.0237 & 0.6225 $\pm$ 0.0416 & 0.6956 $\pm$ 0.0450\\
E16 & 0.6376 $\pm$ 0.0175 & 0.5979 $\pm$ 0.0404 & 0.6748 $\pm$ 0.0202\\
E17 & 0.6452 $\pm$ 0.0219 & 0.6154 $\pm$ 0.0259 & 0.6993 $\pm$ 0.0287\\
\bottomrule\end{tabular}
\end{table}


E11, the CodeBERT BCE backbone comparator, scores 0.6397 on Python, 0.5980 on Java, and 0.6805 on Pharo. It is higher than E1 on Python, nearly equal on Java, and lower on Pharo. Comparing E15 with E11 gives R-Drop-associated increments of about 0.82, 2.45, and 1.51 points within the CodeBERT backbone in this cohort. Those matched comparisons are more informative than attributing the entire E15--E1 difference to regularization.

Loss changes do not improve every language. E3 is the leading Pharo screening configuration but slightly below E1 on Java. Weighted BCE and ASL alter precision/recall tradeoffs rather than uniformly raising micro-F1. E4 is identical to E1 in all reported screening metrics because screening defers threshold tuning; the duplicate row is omitted from the tables rather than treated as a separate gain. More complex graph, attention, and layer-wise optimization variants also fail to dominate all simpler alternatives.

\subsection{Final evaluation}
Tables~\ref{tab:pythontest}--\ref{tab:pharotest} contain all available final results, including the five traditional models. They report micro-F1, macro-F1, Hamming loss, and micro-AP. They do not fill in the twelve missing neural configurations per language.
\begin{table}[htbp]
\centering\small
\caption{Python Stage 2 test results. HL is Hamming loss (lower is better).}\label{tab:pythontest}
\begin{tabular}{lrrrr}\toprule
Method & Micro-F1 & Macro-F1 & HL & Micro-AP\\\midrule
E15 & 0.6541 & 0.5965 & 0.1437 & 0.7585\\
E9 & 0.6285 & 0.5634 & 0.1570 & 0.6779\\
E17 & 0.6248 & 0.5669 & 0.1492 & 0.7474\\
E1 & 0.6074 & 0.5547 & 0.1656 & 0.6571\\
E3 & 0.5900 & 0.5496 & 0.1727 & 0.6880\\
Linear SVM & 0.5704 & 0.5295 & 0.1906 & 0.6326\\
SGD & 0.5690 & 0.5268 & 0.1953 & 0.6395\\
Logistic regression & 0.5675 & 0.5244 & 0.1953 & 0.6392\\
Classifier chain & 0.5577 & 0.5275 & 0.1945 & 0.5772\\
Random forest & 0.5421 & 0.5019 & 0.1953 & 0.6408\\
\bottomrule\end{tabular}
\end{table}

\begin{table}[htbp]
\centering\small
\caption{Java Stage 2 test results. HL is Hamming loss (lower is better).}\label{tab:javatest}
\begin{tabular}{lrrrr}\toprule
Method & Micro-F1 & Macro-F1 & HL & Micro-AP\\\midrule
E17 & 0.7034 & 0.7342 & 0.0921 & 0.8040\\
Classifier chain & 0.6964 & 0.7470 & 0.1004 & 0.7503\\
SGD & 0.6950 & 0.7441 & 0.1015 & 0.7751\\
E1 & 0.6947 & 0.6946 & 0.0945 & 0.7580\\
E15 & 0.6914 & 0.7284 & 0.0933 & 0.7986\\
Logistic regression & 0.6871 & 0.7344 & 0.1027 & 0.7773\\
E3 & 0.6839 & 0.7111 & 0.1009 & 0.7910\\
Linear SVM & 0.6838 & 0.7335 & 0.1015 & 0.7412\\
E9 & 0.6762 & 0.7371 & 0.1068 & 0.7465\\
Random forest & 0.6360 & 0.6495 & 0.1216 & 0.7430\\
\bottomrule\end{tabular}
\end{table}

\begin{table}[htbp]
\centering\small
\caption{Pharo Stage 2 test results. HL is Hamming loss (lower is better).}\label{tab:pharotest}
\begin{tabular}{lrrrr}\toprule
Method & Micro-F1 & Macro-F1 & HL & Micro-AP\\\midrule
E1 & 0.7371 & 0.6720 & 0.0828 & 0.7990\\
E9 & 0.7157 & 0.6058 & 0.0909 & 0.7766\\
E17 & 0.7080 & 0.6283 & 0.0917 & 0.8084\\
E15 & 0.7050 & 0.6194 & 0.0917 & 0.8013\\
E3 & 0.6919 & 0.6015 & 0.0990 & 0.7956\\
Random forest & 0.6447 & 0.5607 & 0.1136 & 0.7126\\
SGD & 0.6420 & 0.5788 & 0.1177 & 0.6618\\
Logistic regression & 0.6400 & 0.5653 & 0.1169 & 0.6691\\
Linear SVM & 0.6308 & 0.5601 & 0.1169 & 0.6545\\
Classifier chain & 0.6161 & 0.5197 & 0.1274 & 0.6224\\
\bottomrule\end{tabular}
\end{table}


Among final neural runs, E15 is highest on Python (micro-F1 0.6541, macro-F1 0.5965), E17 on Java (0.7034, 0.7342), and E1 on Pharo (0.7371, 0.6720). Across all available methods, these are also the highest test micro-F1 values. However, Java's classifier chain has macro-F1 0.7470 and micro-F1 0.6964, illustrating a different tradeoff from the neural test leader.

\subsection{Cross-validation versus test ranking}
Figure~\ref{fig:cvtest} compares the five neural configurations having both stages. Python E15 remains first in this subset. Java E17 overtakes E15; Pharo's E1 overtakes its screening leader E3. CV and final test measure different fitted models and evaluation sets: CV averages selected fold-validation scores, whereas final training uses a new training/validation/test partition. Both training composition and checkpoint selection can affect ranking, as can sampling variation. The observed changes do not establish intrinsic model instability.
\reportfigure{cv_test}{Stage 1 means and Stage 2 test micro-F1 for the five evaluated Transformers. Scores are not independent replications and differences must not be read as generalization-gap estimates.}{fig:cvtest}{1}

\textbf{Answer to RQ2.} E15 leads full neural screening for Python and Java, while E3 leads Pharo. Within the available final subset, E15/E17/E1 yield the highest observed test micro-F1 for Python/Java/Pharo. The ranking depends on the language and evaluated subset.

\subsection{Per-label performance: RQ3}
Tables~\ref{tab:pythonlabels} and \ref{tab:javalabels} compare E1 with the highest observed final micro-F1 Transformer in the corresponding language. This retrospective comparator is descriptive. Figure~\ref{fig:pyheat} and the other heatmaps show all five final configurations so that regressions are not hidden by selecting a single winner.
\begin{table}[htbp]
\centering\small
\caption{Python test label performance. The highest observed test micro-F1 Transformer is E15; comparator choice is retrospective and descriptive.}\label{tab:pythonlabels}
\begin{tabular}{lrrrrrrr}\toprule
Label & Support & E1 P & E1 R & E1 F1 & E15 P & E15 R & E15 F1\\\midrule
DevelopmentNotes & 31 & 0.429 & 0.290 & 0.346 & 0.389 & 0.226 & 0.286\\
Expand & 51 & 0.556 & 0.294 & 0.385 & 0.579 & 0.431 & 0.494\\
Parameters & 80 & 0.642 & 0.650 & 0.646 & 0.743 & 0.688 & 0.714\\
Summary & 45 & 0.771 & 0.600 & 0.675 & 0.680 & 0.756 & 0.716\\
Usage & 80 & 0.685 & 0.762 & 0.722 & 0.862 & 0.700 & 0.772\\
\bottomrule\end{tabular}
\end{table}

\begin{table}[htbp]
\centering\small
\caption{Java test label performance. The highest observed test micro-F1 Transformer is E17; comparator choice is retrospective and descriptive.}\label{tab:javalabels}
\begin{tabular}{lrrrrrrr}\toprule
Label & Support & E1 P & E1 R & E1 F1 & E17 P & E17 R & E17 F1\\\midrule
Expand & 63 & 0.756 & 0.540 & 0.630 & 0.676 & 0.365 & 0.474\\
Ownership & 11 & 1.000 & 1.000 & 1.000 & 1.000 & 1.000 & 1.000\\
Pointer & 35 & 0.913 & 0.600 & 0.724 & 0.833 & 0.571 & 0.678\\
deprecation & 12 & 1.000 & 0.833 & 0.909 & 0.917 & 0.917 & 0.917\\
rational & 28 & 0.600 & 0.107 & 0.182 & 0.857 & 0.429 & 0.571\\
summary & 42 & 0.544 & 0.738 & 0.626 & 0.653 & 0.762 & 0.703\\
usage & 91 & 0.791 & 0.791 & 0.791 & 0.760 & 0.835 & 0.796\\
\bottomrule\end{tabular}
\end{table}

\reportfigure{python_heatmap}{Python test F1 by label and final Transformer configuration. Every value is read from stored per-label metrics; supports are given in Table~\ref{tab:pythonlabels}.}{fig:pyheat}{.95}
\reportfigure{java_heatmap}{Java test F1 by label and configuration. High Ownership scores are based on only eleven positive test assignments.}{fig:jaheat}{.95}
\reportfigure{pharo_heatmap}{Pharo test F1 by label and configuration. E1 leads aggregate test micro-F1, while some of its individual label scores are surpassed.}{fig:phheat}{.95}

On Python, E15 raises aggregate micro-F1 from 0.6074 to 0.6541 and macro-F1 from 0.5547 to 0.5965. Yet DevelopmentNotes F1 falls from 0.346 to 0.286 (31 positives), while Expand rises from 0.385 to 0.494 and Parameters from 0.646 to 0.714. Thus even a simultaneous micro/macro improvement is compatible with a minority-label regression.

On Java, E17 improves micro-F1 over E1 from 0.6947 to 0.7034 and macro-F1 from 0.6946 to 0.7342. The rational label rises from 0.182 to 0.571, while Expand falls from 0.630 to 0.474 and Pointer from 0.724 to 0.678. Ownership is 1.000 for both, based on eleven positives. Per-label recall and precision in Table~\ref{tab:javalabels} show whether a score change reflects missed positives or additional predicted positives; no sentence-level explanation is inferred from counts alone.

On Pharo, E1 has higher aggregate scores than E15, but E15 has higher Keyimplementationpoints F1 (0.537 versus 0.512) and Keymessages F1 (0.677 versus 0.667). Classreferences support is only seven; its F1 ranges from 0.308 for E3 to 0.545 for E1 across the five final configurations. Such rare-category variation is practically relevant but imprecisely estimated. Full E1 per-label precision/recall and support for Pharo are supplied in the appendix.

\textbf{Answer to RQ3.} No. Aggregate improvement does not imply improvement for each label. The stored metrics contain explicit counterexamples in Python and Java, and Pharo shows label-level gains within an aggregate-lower configuration. Neither micro- nor macro-F1 should replace per-label assessment when rare information types matter.


\subsection{Learning dynamics and early stopping}\label{sec:learning}
Figure~\ref{fig:mllearning} uses the saved logistic-regression curve files. Each point is a separate refit at a specified maximum L-BFGS iteration count, not a logged epoch or a guaranteed prefix of one long optimization trajectory. On Python, training Hamming loss falls from 0.3009 at the first limit to 0.0178 at the final limit; validation micro-F1 settles near 0.5654. The corresponding final validation values are 0.6046 for Java and 0.6294 for Pharo. Strong training fit with a lower validation plateau indicates limited generalization benefit from further optimization; the different metrics must not be subtracted as if they were the same loss.
\reportfigure{ml_learning}{Stored logistic-regression refit diagnostics. Top: training Hamming loss; bottom: validation F1. The horizontal axis is the maximum solver iteration count on a log scale, not a neural epoch.}{fig:mllearning}{1}

Figure~\ref{fig:neurallearning} shows the actual E1/E15 training histories. For Java E15, validation micro-F1 peaks at epoch 3 (0.6402) and declines through epochs 4--5 while training objective continues falling; patience 2 stops the run after epoch 5 and restores epoch 3. Java E1 similarly selects epoch 6 and finishes at epoch 8. Python E1/E15 both select epoch 9 out of 10; Pharo E1 selects epoch 8 out of 10 and E15 epoch 6 out of 8. These trajectories demonstrate why final-epoch weights are not necessarily the selected weights. A falling training objective with weaker validation scores is consistent with overfitting, though one split does not establish a universal instability property. E1 and E15 have different objectives, so their training-loss magnitudes are not directly comparable.
\reportfigure{neural_learning}{Saved neural histories for E1 and E15. Stars mark selected validation epochs. Loss values are objective-specific; the bottom row compares the common validation micro-F1 metric.}{fig:neurallearning}{1}

\subsection{Illustrative comments and quantitative difficult-label analysis}\label{sec:illustrations}
Table~\ref{tab:examples} gives actual dataset sentences and their annotated labels. Python sentence 19 has both Expand and Usage; Java sentence 61 combines rational and summary. Such examples make the multi-label formulation concrete and show that annotation can depend on fragments and surrounding context. They are not claimed to be model errors or model-generated explanations.
\begin{table}[htbp]\centering\small
\caption{Actual annotated sentences illustrating information types. These are dataset examples, not model errors or predictions.}\label{tab:examples}
\begin{tabularx}{\linewidth}{llXX}\toprule Language & ID & Sentence & Annotated labels\\\midrule
Python & 8 & only provides a label and a dict of models. & Parameters\\
Python & 19 & decorated function in the instance itself, used as a context manager. & Expand, Usage\\
Java & 61 & contains all the state data that needs to be stored persistently for an application & rational, summary\\
Pharo & 3 & i contain examples of different types of blarrowheadelement & Classreferences, Responsibilities\\
\bottomrule\end{tabularx}\end{table}


The Java rational label provides a quantitative difficult case. E1 detects 3 of 28 positives (recall 0.1071), predicts 5 positives, and therefore has 25 false negatives and 2 false positives. E17 detects 12 of 28 (recall 0.4286), predicts 14 positives, and has 16 false negatives and 2 false positives. Its F1 rises from 0.1818 to 0.5714 through more true positives without an increase in false-positive count on this split. This is a count-level explanation supported by stored confusion evidence; the missing row-level predictions prevent identifying which sentences changed.

\subsection{Published benchmark context}\label{sec:published}
The organizers' values in Table~\ref{tab:official} average 19 binary category F1 scores on official partitions~\cite{nlbse_overview}. Their ranking weights average F1 by 0.75 and the proportion of improved categories by 0.25. Our per-language micro-F1 uses different splits and aggregation, so these published values locate the task without supplying a directly comparable superiority margin.
\begin{table}[htbp]\centering\small
\caption{Published NLBSE'23 results from the organizers' Table III; context only, not directly comparable with our grouped splits.}\label{tab:official}
\begin{tabular}{lrrr}\toprule Method & Average precision & Average recall & Average F1\\\midrule
Official RF baseline & 0.439 & 0.245 & 0.309\\
Submitted logistic regression & 0.540 & 0.560 & 0.547\\
CodeT5 & 0.728 & 0.606 & 0.657\\
STACC & 0.710 & 0.794 & 0.744\\\bottomrule
\end{tabular}\end{table}
